\subsection{The eight-state P5--P5 scope}

\begin{lemma}[Proved: a unique regular P5 core opposite P5]
\label{CP2:A-P5-double-unique-core}
Let $S$ be any finite nonempty left-$P_5$, right-$P_5$
biset. Suppose the right action has no fixed state and its
minimal-band core is a unique regular two-point column
orbit. The collage cannot have rational projective-plane
cohomology.
\end{lemma}
\begin{proof}
Write $P=P_5$, $E=B(P,S,P)$ and $X$ for the collage.
The semilinear resolution computation in
Lemma~\ref{CP2:A-mixed-five-reductions}(3) is equivariant
for any commuting opposite action. Before taking its
derived coinvariants it gives the complex
$\mathbb Q[0]\oplus T[1]$, with signs $+1,-1$ under
the source-row interchange $F$, respectively. Therefore
\[
 H_3(E;\mathbb Q)=H_3(BP;\mathbb Q)\oplus
 \operatorname{Tor}^{\mathbb QP}_2(\mathbb Q,T).
\]
The first summand is zero, so $F_*$ is $-1$ on this group.
Both corner spaces have zero rational homology in degrees
three and four by their length-two projective resolution.
The mapping-cylinder exact sequence consequently identifies
$H_4(X;\mathbb Q)$ with $H_3(E;\mathbb Q)$, with sign
$-1$. If $H^2(X;\mathbb Q)$ is one-dimensional, its
generator is an eigenvector of the involution $F^*$ with
eigenvalue $+1$ or $-1$. Its square has eigenvalue $+1$
in either case and must be zero. This contradicts the
projective-plane cohomology ring. No acyclicity assumption
on the opposite copy of $BP$ was used.
\end{proof}

\begin{proposition}[Proved: the complete finite right-act scope at $(5,5,8)$]
\label{CP2:A-P5-double-eight-scope}
To exclude the P5--P5 portion of the eighteen-arrow profile
$(5,5,8)$, it suffices to enumerate every left $P_5$ action
commuting with each of the twenty-nine specified eight-state
right acts. They are the right normal forms indexed
$1,2,\ldots,29$ in the checked eight-state catalog
\path{monoid/certificates/p5_small_targets/right_acts_8.tsv}.
State automorphisms may be divided out, with exact orbit
weights. No corner automorphism or opposition quotient is
needed in that finite enumeration.
\end{proposition}
\begin{proof}
If both actions have fixed states, the source and target
corner inclusions are separately nullhomotopic in the
collage, by the corresponding natural transformations.
Every second class therefore restricts to zero on both
corners and has zero square by the relative-cup lemma.
Otherwise at least one action has no fixed state.
Opposition and $P_5\cong P_5^{\mathrm{op}}$ put such an
action on the right, preserving the classifying-space
homotopy type. This step does not assert that both actions
are fixed-point free.

The independently checked eight-state normal-form catalog
has $1212$ ordinary right-act classes, of which exactly
the first $31$ have no fixed state. Index zero has a
unique two-point core and is excluded by the preceding
lemma. Index thirty consists of four regular two-point
cores and no transient states, so the full-core wedge
lemma applies. The other twenty-nine forms are the stated
scope. The cost inventory
\path{monoid/certificates/p5_double18_cost_inventory.json}
reconstructs their core types and counts every commuting
endomorphism from its core-orbit map and allowed transient
images; their largest centralizer has $2624$ elements.
The normal-form and Burnside receipt remains the earlier
complete checked eight-state receipt. A completed finite
search, its independent scope check and its literal cup
certificates are separate obligations.
\end{proof}
